\documentclass[11pt, a4paper]{article}

\usepackage[utf8]{inputenc}
\usepackage[T1]{fontenc}
\usepackage[spanish]{babel}
\usepackage{lmodern}
\usepackage[margin=2.2cm]{geometry}
\usepackage{xcolor}
\usepackage{fancyhdr}
\usepackage{enumitem}
\usepackage{titlesec}
\usepackage[most]{tcolorbox}
\usepackage{amssymb, amsmath}
\usepackage{parskip}
\usepackage{hyperref}
\usepackage{needspace}
\usepackage{booktabs}
\usepackage{array}

% ---- Paleta ----
\definecolor{verdeTerm}{HTML}{1C8A3F}
\definecolor{verdeOsc}{HTML}{146A2F}
\definecolor{tinta}{HTML}{0B0B0B}
\definecolor{tinta2}{HTML}{4A4948}
\definecolor{mute}{HTML}{8F8D87}
\definecolor{plano}{HTML}{F7F4EB}
\definecolor{planoOsc}{HTML}{EFEBDD}
\definecolor{linea}{HTML}{DAD6C8}
\definecolor{ambar}{HTML}{C78A1E}
\definecolor{rojoSoft}{HTML}{B23A48}

\hypersetup{
  colorlinks=true, linkcolor=verdeTerm, urlcolor=verdeTerm,
  pdftitle={Manual teorico - Defensa TFM data.lung},
  pdfauthor={Daniel Tapia Diez}
}

\titleformat{\section}
  {\Large\bfseries\color{tinta}}
  {\thesection}{0.5em}{}
  [\vspace{-0.4em}{\color{verdeTerm}\rule{2.5cm}{2pt}}]

\titleformat{\subsection}
  {\large\bfseries\color{verdeOsc}}
  {\thesubsection}{0.5em}{}

\titleformat{\subsubsection}
  {\normalsize\bfseries\color{tinta}}
  {\thesubsubsection}{0.5em}{}

\pagestyle{fancy}
\fancyhf{}
\fancyhead[L]{\small\ttfamily\color{verdeTerm}\textbf{>}\color{tinta2}\ data\textcolor{verdeTerm}{.}lung}
\fancyhead[R]{\small\color{tinta2}Manual te\'orico $\cdot$ defensa TFM}
\fancyfoot[C]{\small\color{mute}\thepage}
\renewcommand{\headrulewidth}{0.4pt}
\renewcommand{\footrulewidth}{0pt}
\setlength{\headheight}{15pt}

\newtcolorbox{concepto}[1]{
  enhanced, breakable,
  colback=plano, colframe=verdeTerm,
  boxrule=1pt, arc=3pt,
  left=8pt, right=8pt, top=6pt, bottom=6pt,
  title={\bfseries\color{white}#1},
  colbacktitle=verdeTerm,
  attach boxed title to top left={xshift=6pt, yshift=-6pt},
  boxed title style={colback=verdeTerm, sharp corners, boxrule=0pt}
}
\newtcolorbox{peligro}{enhanced, breakable,
  colback=planoOsc, colframe=rojoSoft,
  boxrule=0.9pt, arc=2pt,
  left=8pt, right=8pt, top=5pt, bottom=5pt}
\newtcolorbox{tipbox}{enhanced, breakable,
  colback=planoOsc, colframe=ambar,
  boxrule=0.8pt, arc=2pt,
  left=8pt, right=8pt, top=5pt, bottom=5pt}
\newtcolorbox{formula}{enhanced, breakable,
  colback=white, colframe=linea,
  boxrule=0.6pt, arc=2pt,
  left=8pt, right=8pt, top=6pt, bottom=6pt}

\newcommand{\gene}[1]{\textit{#1}}
\newcommand{\slidebadge}[1]{{\colorbox{verdeTerm}{\color{white}\bfseries\;#1\;}}}

\begin{document}

% =================== PORTADA =========================
\thispagestyle{empty}
\vspace*{1.5cm}
\begin{center}
  {\ttfamily\fontsize{44pt}{44pt}\selectfont\bfseries
   data\textcolor{verdeTerm}{.}lung}\\[0.2cm]
  {\itshape\normalsize\color{tinta2}Del dataset a una firma g\'enica replicable}\\[0.8cm]

  \textcolor{verdeTerm}{\rule{3cm}{2pt}}\\[0.6cm]

  {\LARGE\bfseries Manual te\'orico}\\[0.3cm]
  {\large\color{tinta2}Todo lo que necesitas saber\\para defenderte de cualquier pregunta}\\[1.5cm]

  \begin{tabular}{rl}
    \color{mute}\textbf{Autor}      & Daniel Tapia D\'iez \\
    \color{mute}\textbf{Universidad}& UAX $\cdot$ M\'aster en Bioinform\'atica \\
    \color{mute}\textbf{Defensa}    & 14 de agosto de 2026 \\
  \end{tabular}\\[2cm]

  \begin{tcolorbox}[colback=plano, colframe=verdeTerm, width=0.85\textwidth,
                    arc=3pt, boxrule=1pt]
    \small\itshape
    Este documento cubre los fundamentos te\'oricos de biolog\'ia,
    tecnolog\'ia, estad\'istica, aprendizaje autom\'atico y modelos de
    lenguaje que soportan el TFM. Se organiza en 12 partes tem\'aticas,
    cada una autocontenida. \'Usalo como referencia r\'apida ante
    cualquier pregunta del tribunal.
  \end{tcolorbox}
\end{center}
\newpage

\tableofcontents
\newpage

% ========================================================
\section{C\'ancer de pulm\'on: contexto cl\'inico}

\subsection{Epidemiolog\'ia}

El c\'ancer de pulm\'on es la primera causa de muerte por c\'ancer en el
mundo. En 2022 se diagnosticaron $\approx 2{,}48$ millones de casos
(GLOBOCAN, IARC) con $\approx 1{,}82$ millones de fallecimientos anuales
atribuibles a la enfermedad. La supervivencia global a 5 a\~nos apenas
supera el 20\,\%, muy por debajo de otros tumores s\'olidos (mama
$\approx 90\,\%$, pr\'ostata $\approx 98\,\%$, colon $\approx 65\,\%$).

Factores de riesgo principales: tabaco (85\,\% de los casos), exposici\'on
a rad\'on, amianto, contaminaci\'on atmosf\'erica, historia familiar y
alteraciones gen\'eticas hereditarias (aunque el componente hereditario
es minoritario). En no fumadores, especialmente en mujeres asi\'aticas,
predominan mutaciones activadoras en EGFR y ALK.

\subsection{Clasificaci\'on histol\'ogica}

El c\'ancer de pulm\'on se divide en dos grandes grupos:

\begin{itemize}[leftmargin=1.2em]
  \item \textbf{C\'ancer de pulm\'on de c\'elulas peque\~nas} (SCLC,
  \emph{Small Cell Lung Cancer}). Aproximadamente el 15\,\% de los casos.
  Muy agresivo, altamente asociado al tabaco, con tratamiento
  quimioter\'apico distinto y peor pron\'ostico.
  \item \textbf{C\'ancer de pulm\'on no microc\'itico} (NSCLC,
  \emph{Non-Small Cell Lung Cancer}). El 85\,\% restante. Se subdivide en:
  \begin{itemize}
    \item \textbf{Adenocarcinoma (ADC)}: 40--50\,\% de los NSCLC. Origen
    glandular. Aparece en la periferia del pulm\'on. M\'as com\'un en no
    fumadores y en mujeres. Marcadores IHC: TTF-1/NKX2-1, napsina A.
    \item \textbf{Carcinoma escamoso (SQC, LUSC)}: 25--30\,\%. Origen en
    el epitelio bronquial. M\'as central. Fuerte asociaci\'on con el
    tabaco. Marcadores IHC: p40 (isoforma TP63), CK5/6, DSG3.
    \item \textbf{Carcinoma de c\'elulas grandes}: pooco frecuente,
    diagn\'ostico de exclusi\'on.
    \item \textbf{Tumores neuroendocrinos}: carcinoide t\'ipico y
    at\'ipico. Perfil transcript\'omico muy distinto.
  \end{itemize}
\end{itemize}

\begin{concepto}{Por qu\'e distinguir ADC de SQC es cr\'itico}
El pemetrexed (antifolato) est\'a contraindicado en escamoso porque no
funciona con el mismo mecanismo bioqu\'imico. El bevacizumab
(antiangiog\'enico anti-VEGF) est\'a contraindicado en escamoso por
\textbf{riesgo de hemorragia pulmonar mortal}. Diagnosticar
err\'oneamente un escamoso como adenocarcinoma puede matar al paciente
en las primeras semanas de tratamiento. Por eso todo el pipeline se
centra en esta separaci\'on binaria.
\end{concepto}

\subsection{Diagn\'ostico cl\'inico actual}

El est\'andar cl\'inico es la biopsia con estudio anatomopatol\'ogico e
inmunohistoqu\'imica (IHC). La OMS publica desde 2015 un panel de
marcadores recomendados (Travis et al., 2015):

\begin{itemize}[leftmargin=1.2em]
  \item \textbf{Escamoso}: KRT5, KRT6A, KRT6B, KRT13, KRT14, KRT15, KRT16,
  KRT17, CDH3, DSG3, PKP1, TP63 (12 marcadores).
  \item \textbf{Adenocarcinoma}: NKX2-1 (TTF-1), NAPSA (napsina A), SFTPB,
  SFTPC, SFTPA1, SFTPA2, MUC1, CEACAM5, CEACAM6, KRT7 (8-10 marcadores).
\end{itemize}

Estos marcadores se detectan mediante anticuerpos con revelado
crom\'ogeno. El pat\'ologo mira al microscopio la muestra te\~nida y
determina el subtipo bas\'andose en el patr\'on de tinci\'on.

\subsection{Estadificaci\'on TNM}

Sistema TNM (\emph{Tumor, Nodes, Metastasis}) de la AJCC. T describe
tama\~no e invasi\'on local, N el compromiso ganglionar, M la presencia
de met\'astasis a distancia. Estadios I--IV. En este TFM las cohortes
mezclan estadios variables y se trabaja sobre la clasificaci\'on
histol\'ogica, no la TNM.

% ========================================================
\section{Biolog\'ia molecular esencial}

\subsection{ADN, ARN, prote\'ina}

El \textbf{dogma central de la biolog\'ia molecular} (Crick, 1958)
establece el flujo de informaci\'on:

\[ \text{ADN} \xrightarrow{\text{transcripci\'on}} \text{ARN}
   \xrightarrow{\text{traducci\'on}} \text{prote\'ina} \]

\begin{itemize}[leftmargin=1.2em]
  \item \textbf{ADN} (\'acido desoxirribonucleico): mol\'ecula que
  contiene el genoma. Doble h\'elice de nucle\'otidos. En humanos, ${\sim}3$
  mil millones de pares de bases y ${\sim}20\,000$ genes codificantes de
  prote\'ina.
  \item \textbf{Gen}: fragmento de ADN que codifica una prote\'ina o una
  funci\'on reguladora (por ejemplo un ARN no codificante).
  \item \textbf{ARN mensajero (mRNA)}: copia de un gen transcrito, que
  sale del n\'ucleo hacia el citoplasma para ser traducido a prote\'ina.
  \item \textbf{Prote\'ina}: producto funcional. Ejecuta la mayor parte
  de las funciones celulares (estructural, cataltica, se\~nalizaci\'on).
\end{itemize}

\subsection{Expresi\'on g\'enica y transcriptoma}

Aunque todas las c\'elulas de un individuo comparten el mismo ADN, no
transcriben los mismos genes al mismo tiempo. La c\'elula \emph{expresa}
un gen cuando lo transcribe activamente a mRNA. El conjunto de todos los
mRNAs presentes en una muestra biol\'ogica en un momento dado se llama
\textbf{transcriptoma}. Es un retrato din\'amico del estado funcional de
la c\'elula.

Medir el transcriptoma equivale a medir cu\'antas copias de mRNA hay de
cada gen. Los valores se expresan t\'ipicamente en escala logar\'itmica
($\log_2$) porque abarcan varios \'ordenes de magnitud y porque en escala
log las diferencias biol\'ogicas se aproximan a una distribuci\'on
gaussiana.

\begin{concepto}{Qu\'e ``mide'' realmente el transcriptoma}
No mide directamente actividad celular. Mide \textbf{abundancia relativa
de mRNA}. Hay una correlaci\'on razonable entre nivel de mRNA y nivel de
prote\'ina, pero no perfecta: los mRNAs pueden estar retenidos, no
traducirse, o las prote\'inas pueden ser muy estables o degradarse
r\'apidamente. En el TFM asumimos --- como toda la literatura de firmas
g\'enicas --- que las diferencias sostenidas de mRNA reflejan cambios
funcionales relevantes.
\end{concepto}

\subsection{S\'imbolos g\'enicos oficiales}

Cada gen humano tiene un s\'imbolo oficial asignado por HGNC (\emph{HUGO
Gene Nomenclature Committee}). Ejemplos: \gene{TP53}, \gene{EGFR},
\gene{KRT5}, \gene{DSG3}. Los s\'imbolos evolucionan con el tiempo (algunos
se renombran); las bases de datos de expresi\'on usan la versi\'on
vigente en el momento en que el fabricante del chip anot\'o las sondas.

% ========================================================
\section{Tecnolog\'ias de medici\'on de expresi\'on}

\subsection{Microarrays}

Un microarray es un chip de vidrio del tama\~no de una diapositiva con
decenas de miles de \textbf{sondas} de ADN adheridas en posiciones
conocidas. Cada sonda es una secuencia de $\sim$25 nucle\'otidos
complementaria a un fragmento de un mRNA espec\'ifico. Al hibridar una
muestra biol\'ogica marcada con fluorescencia:

\begin{enumerate}[leftmargin=1.4em]
  \item Se extrae ARN total de la muestra.
  \item Se retrotranscribe a cDNA con fluoref\'oros.
  \item Se hibrida sobre el chip: los cDNAs se unen a sus sondas
  complementarias.
  \item Se escanea la fluorescencia: cada punto brilla proporcionalmente
  a la cantidad de cDNA que se uni\'o.
\end{enumerate}

Plataforma m\'as com\'un en el TFM: \textbf{Affymetrix HG-U133 Plus 2.0}
(identificador GEO \texttt{GPL570}), $\sim$54\,675 sondas para ${\sim}20\,000$
genes. Redundancia (varias sondas por gen) para robustez estad\'istica.

\subsubsection*{Ventajas y limitaciones}
Ventajas: barato, r\'apido, gran base hist\'orica de datos.
Limitaciones: rango din\'amico comprimido, dif\'icil detectar
transcritos nuevos, sesgo por secuencia de sonda.

\subsection{RNA-Seq}

\textbf{RNA-Seq} secuencia directamente el mRNA:

\begin{enumerate}[leftmargin=1.4em]
  \item Se extrae mRNA y se fragmenta.
  \item Se retrotranscribe a cDNA.
  \item Se secuencian los fragmentos generando \emph{reads} de $\sim$100
  bases (Illumina).
  \item Cada read se alinea al genoma o transcriptoma de referencia.
  \item Se cuentan reads por gen (o por transcrito).
\end{enumerate}

Ventajas: alto rango din\'amico, detecta variantes nuevas, cuantitativo.
Limitaciones: coste (aunque cayendo), requerimientos computacionales.

\textbf{Unidades t\'ipicas}:
\begin{itemize}[leftmargin=1.2em]
  \item \textbf{Raw counts}: n\'umero de reads por gen.
  \item \textbf{TPM} (\emph{Transcripts Per Million}): normaliza por longitud del gen y por profundidad de secuenciaci\'on.
  \item \textbf{FPKM} (\emph{Fragments Per Kilobase per Million}): similar, para paired-end.
  \item \textbf{RSEM}: estimador estad\'istico que reparte reads que mapean ambiguamente. Usado por TCGA.
\end{itemize}

En el TFM: entrenamos en microarray, validamos en RNA-Seq (TCGA con RSEM).
Que un panel funcione en ambas plataformas es evidencia de que la se\~nal
es \textbf{biol\'ogica y no artefactual}.

% ========================================================
\section{GEO, TCGA y datos p\'ublicos}

\subsection{Gene Expression Omnibus (GEO)}

Repositorio del NCBI (EE.\ UU.) que aloja m\'as de 200\,000 estudios de
expresi\'on g\'enica desde el a\~no 2000. Terminolog\'ia:

\begin{itemize}[leftmargin=1.2em]
  \item \textbf{GSE} (\emph{Gene Series}): un estudio completo, ejemplo
  \texttt{GSE31210}.
  \item \textbf{GSM} (\emph{Gene Sample}): una muestra individual dentro
  de un estudio.
  \item \textbf{GPL} (\emph{Gene Platform}): la plataforma tecnol\'ogica
  usada (chip Affymetrix, etc.). Cada GPL trae su tabla de anotaci\'on
  con la correspondencia sonda $\to$ gen.
\end{itemize}

Los metadatos cl\'inicos en GEO est\'an en \textbf{texto libre} sin
estructura est\'andar. Cada laboratorio los redacta como quiere: en
\texttt{characteristics\_ch1} puede aparecer ``\emph{primary lung tumor,
LUSC, stage III}'' o ``\emph{sample from patient with squamous cell
carcinoma}''. Automatizar la clasificaci\'on requiere leer y entender
lenguaje natural: por eso el TFM usa un LLM.

\subsection{The Cancer Genome Atlas (TCGA)}

Gran consorcio del NIH americano que caracteriz\'o molecularmente ${\sim}11\,000$
tumores humanos entre 2006 y 2018. Para pulm\'on hay dos proyectos:

\begin{itemize}[leftmargin=1.2em]
  \item \textbf{TCGA-LUAD}: $\sim$517 adenocarcinomas con expresi\'on
  RNA-Seq, mutaciones, metilaci\'on, cl\'inicos.
  \item \textbf{TCGA-LUSC}: $\sim$504 escamosos, misma caracterizaci\'on.
\end{itemize}

En el TFM se usan 230 LUAD + 178 LUSC = 408 muestras a trav\'es de la
\textbf{API REST de cBioPortal}, que sirve datos TCGA ya limpios y
homogeneizados. La transferencia del panel LASSO entrenado en microarray
a TCGA RNA-Seq da AUC = 0,857 con IC 95\,\% [0,812; 0,899].

% ========================================================
\section{Preprocesado de expresi\'on: normalizaci\'on}

Antes de cualquier an\'alisis estad\'istico hay tres transformaciones
est\'andar sobre matrices de expresi\'on.

\subsection{Transformaci\'on logar\'itmica}

Las intensidades de sonda de microarray abarcan varios \'ordenes de
magnitud (10 a 65\,000 en escala cruda). Aplicamos $x' = \log_2(x+1)$:

\begin{itemize}[leftmargin=1.2em]
  \item Aproxima las distribuciones a gaussianas (requisito de tests
  param\'etricos).
  \item Convierte diferencias multiplicativas en aditivas: una diferencia
  de $2 \times$ en escala cruda equivale a $+1$ en escala $\log_2$.
  \item El ``+1'' evita $\log_2(0) = -\infty$.
\end{itemize}

En RNA-Seq (counts), el est\'andar equivalente es $\log_2(TPM+1)$ o
transformaciones variance-stabilizing (VST, DESeq2).

\subsection{Normalizaci\'on por cuantiles}

Iguala las distribuciones de intensidad entre muestras. Algoritmo:

\begin{enumerate}[leftmargin=1.4em]
  \item Ordenar cada muestra de menor a mayor intensidad.
  \item Reemplazar cada valor por la media de esa posici\'on entre todas
  las muestras.
  \item Devolver los valores a la posici\'on original de cada muestra.
\end{enumerate}

Resultado: todas las muestras tienen exactamente la misma distribuci\'on
(mismo quantil, misma media, misma mediana). Asume que la distribuci\'on
\emph{global} de expresi\'on es comparable entre muestras --- razonable
cuando la mayor\'ia de genes no cambia entre condiciones.

Alternativas: RMA (Robust Multi-array Average, espec\'ifica de
Affymetrix), MAS5, VSN. En el TFM usamos $\log_2$ + cuantiles por su
simplicidad y aceptaci\'on hist\'orica.

\subsection{Por qu\'e normalizar dentro de cada estudio y no entre estudios}

Normalizar por cuantiles \emph{entre estudios} es una forma implicita de
correcci\'on de lote: fuerza a que las cohortes tengan la misma
distribuci\'on, borrando diferencias t\'ecnicas... y biol\'ogicas. Si
una cohorte est\'a enriquecida en tumores agresivos, cuantilar contra
otra menos agresiva puede borrar la se\~nal. Por eso normalizamos
\emph{dentro} de cada cohorte y confiamos en la exigencia del consenso
multi-cohorte para descartar artefactos.

% ========================================================
\section{Estad\'istica para expresi\'on g\'enica}

\subsection{Pruebas de hip\'otesis: Welch, Mann-Whitney}

Comparar la expresi\'on media de un gen entre dos grupos (tumor vs sano):

\begin{itemize}[leftmargin=1.2em]
  \item \textbf{t de Student} asume varianzas iguales y normalidad. Poco
  robusto.
  \item \textbf{t de Welch}: modificaci\'on que no asume varianzas
  iguales. Es el est\'andar en expresi\'on g\'enica log-transformada.
  Usado en el paso 5 del pipeline.
  \item \textbf{Mann-Whitney U}: alternativa no param\'etrica. \'Util
  cuando la distribuci\'on es muy asim\'etrica.
\end{itemize}

\begin{formula}
\textbf{Estad\'istico $t$ de Welch}:
\[ t = \frac{\bar{x}_1 - \bar{x}_2}{\sqrt{\dfrac{s_1^2}{n_1} + \dfrac{s_2^2}{n_2}}} \]

Los grados de libertad se calculan con la aproximaci\'on de
Welch-Satterthwaite:
\[ \nu \approx \frac{\left(\frac{s_1^2}{n_1} + \frac{s_2^2}{n_2}\right)^2}{\dfrac{(s_1^2/n_1)^2}{n_1-1} + \dfrac{(s_2^2/n_2)^2}{n_2-1}} \]
\end{formula}

\textbf{Cu\'ando cada uno}:

\begin{itemize}[leftmargin=1.2em]
  \item Con $n_i \geq 30$ por grupo y datos log-transformados: t de Welch
  robusto.
  \item Con $n_i < 20$ o distribuciones muy sesgadas: Mann-Whitney U
  m\'as conservador.
  \item Si se detecta bimodalidad severa: pruebas espec\'ificas de mezcla
  (ejemplo: modelos de mezcla gaussiana).
\end{itemize}

\begin{peligro}
Nunca se debe seleccionar el test tras mirar los datos y ver cu\'al da
mejor p-valor: eso es \emph{p-hacking}. El TFM prefija Welch como test
principal para toda la matriz de expresi\'on, con FDR de Benjamini-Hochberg.
\end{peligro}

\subsubsection*{Supuestos, robustez y qu\'e pasa si se violan}

Los supuestos de la $t$ de Welch son:
\begin{enumerate}[leftmargin=1.4em]
  \item Muestras independientes entre y dentro de grupos.
  \item Distribuci\'on aproximadamente sim\'etrica (no necesariamente
  gaussiana con $n$ grande, por el TLC).
  \item Sin outliers extremos.
\end{enumerate}

Violaciones y consecuencias:
\begin{itemize}[leftmargin=1.2em]
  \item \textbf{Dependencia entre muestras} (ejemplo: r\'eplicas t\'ecnicas
  del mismo paciente): inflan la significancia. Se soluciona con modelos
  mixtos (\texttt{lmer}) o promediando r\'eplicas antes del test.
  \item \textbf{Asimetr\'ia severa}: el t-test sigue siendo v\'alido con
  $n \geq 30$ por grupo pero pierde potencia. Alternativa: Mann-Whitney.
  \item \textbf{Outliers}: el t-test es sensible. En gen\'omica, un
  outlier de expresi\'on suele ser artefacto t\'ecnico. La normalizaci\'on
  por cuantiles y el $\log_2$ los reducen; queda como l\'imite conocido.
\end{itemize}

\subsection{Correcci\'on de m\'ultiples pruebas: FDR}

Cuando se hacen 20\,000 tests (uno por gen), a $\alpha = 0{,}05$ se
esperan 1\,000 falsos positivos por azar. Soluciones:

\begin{itemize}[leftmargin=1.2em]
  \item \textbf{Bonferroni}: dividir $\alpha$ por el n\'umero de tests.
  Muy conservador, elimina se\~nal real.
  \item \textbf{Benjamini-Hochberg (BH-FDR)}: controla la
  \emph{proporci\'on esperada de falsos positivos entre los rechazos}.
  Menos conservador, est\'andar en gen\'omica. El TFM usa BH-FDR.
\end{itemize}

\begin{formula}
El \emph{q-value} (FDR-adjusted) de un gen se calcula ordenando los
p-valores de menor a mayor y multiplicando cada uno por $m/i$ (con $m$ =
n\'umero total de tests, $i$ = posici\'on en el ranking). Un q-value $<
0{,}05$ significa: ``en el conjunto de genes con q $< 0{,}05$, esperamos
menos del 5\,\% de falsos positivos''.
\end{formula}

\subsubsection*{FWER vs FDR}

\begin{itemize}[leftmargin=1.2em]
  \item \textbf{FWER} (Family-Wise Error Rate): probabilidad de cometer
  \emph{al menos un} falso positivo. Controlado por Bonferroni.
  \item \textbf{FDR} (False Discovery Rate): proporci\'on \emph{esperada}
  de falsos positivos entre los resultados declarados positivos.
  Controlado por BH.
\end{itemize}

\textbf{Cu\'ando cada uno}:
\begin{itemize}[leftmargin=1.2em]
  \item FWER cuando el coste de un falso positivo es cr\'itico (ensayos
  cl\'inicos regulatorios).
  \item FDR cuando se generan hip\'otesis para investigaci\'on posterior
  (an\'alisis exploratorio de expresi\'on). \textbf{Est\'andar en
  gen\'omica}.
\end{itemize}

Otros m\'etodos: Holm-Bonferroni (secuencial FWER, m\'as potente que
Bonferroni), Storey-Tibshirani $q$-value (versi\'on m\'as sofisticada
de FDR con estimaci\'on de $\pi_0$).

\subsection{Tama\~no de efecto: $d$ de Cohen}

El p-valor solo dice ``la diferencia es estad\'isticamente significativa''.
No dice si la diferencia es \textbf{importante}. La $d$ de Cohen mide
el tama\~no de efecto en unidades de desviaci\'on t\'ipica:

\begin{formula}
\[ d = \frac{\bar{x}_{\text{grupo A}} - \bar{x}_{\text{grupo B}}}{s_{\text{pooled}}} \]
donde $s_{\text{pooled}}$ es la desviaci\'on t\'ipica combinada.
Convenciones: $|d| < 0{,}2$ trivial, $0{,}2$-$0{,}5$ peque\~no, $0{,}5$-$0{,}8$
moderado, $> 0{,}8$ grande.
\end{formula}

En el TFM: un gen entra en la firma de consenso s\'olo si $|d| > 0{,}5$
y con el mismo signo en las 3 cohortes independientes.

\subsubsection*{Otras medidas de tama\~no de efecto}

\begin{itemize}[leftmargin=1.2em]
  \item \textbf{Hedges $g$}: variante de Cohen $d$ corregida por sesgo
  para muestras peque\~nas. Multiplica $d$ por $J = 1 - 3/(4(n_1+n_2)-9)$.
  \item \textbf{Glass $\Delta$}: usa solo la desviaci\'on del grupo
  control. \'Util cuando los grupos difieren mucho en dispersi\'on.
  \item \textbf{Log$_2$ Fold Change (log2FC)}: est\'andar en RNA-Seq.
  Simplemente la diferencia de medias en escala log. Equivalente al
  numerador de $d$ pero sin estandarizar.
\end{itemize}

\begin{concepto}{Por qu\'e $d$ y no log2FC en el TFM}
El log2FC es la cantidad ``biol\'ogica'' (magnitud del cambio) pero no
tiene en cuenta la dispersi\'on. Dos genes pueden tener el mismo log2FC
pero uno con dispersi\'on pequenya (efecto claro) y otro con dispersi\'on
enorme (mero ruido). El $d$ de Cohen combina magnitud y dispersi\'on en
un solo n\'umero, m\'as robusto para meta-an\'alisis por consenso.
\end{concepto}

\subsection{Correlaci\'on de Spearman ($\rho$)}

Mide asociaci\'on mon\'otona entre dos variables ordinales:

\[ \rho = 1 - \frac{6 \sum d_i^2}{n(n^2-1)} \]

donde $d_i$ es la diferencia de rangos. Rango $[-1, +1]$. Diferencia
crucial vs Pearson: Spearman no asume linealidad, solo monotonicidad.
\'Util cuando las escalas no son directamente comparables.

Uso en el TFM: correlacionar el score LASSO con marcadores de composici\'on
alv\'eolo-capilar. $\rho = -0{,}626$ media (4 cohortes con muestra
suficiente), hasta $\rho = -0{,}858$ en GSE31210 con $p = 1{,}1 \cdot 10^{-66}$.

\subsubsection*{Spearman vs Pearson}

\begin{itemize}[leftmargin=1.2em]
  \item \textbf{Pearson $r$}: asume relaci\'on \emph{lineal}. Sensible a
  outliers. \'Util cuando las dos variables est\'an en escalas
  comparables y la relaci\'on es lineal.
  \item \textbf{Spearman $\rho$}: transforma a rangos y aplica Pearson
  sobre rangos. Captura cualquier relaci\'on \emph{mon\'otona}, no s\'olo
  lineal. Robusto frente a outliers. \textbf{Est\'andar en gen\'omica}.
  \item \textbf{Kendall $\tau$}: alternativa a Spearman m\'as robusta,
  computacionalmente m\'as cara.
\end{itemize}

\subsection{Intervalos de confianza y bootstrap}

Un \textbf{IC 95\,\%} no significa ``hay 95\,\% de probabilidad de que el
valor real est\'e en el intervalo''. Significa: ``si repiti\'eramos el
experimento muchas veces, el 95\,\% de los intervalos resultantes
contendr\'ian el valor verdadero''.

\textbf{Bootstrap}: t\'ecnica de remuestreo para estimar la distribuci\'on
de una m\'etrica sin asumir distribuci\'on te\'orica:

\begin{enumerate}[leftmargin=1.4em]
  \item Del conjunto original de $n$ muestras, se toma una muestra de
  $n$ con reemplazo (bootstrap sample).
  \item Se calcula la m\'etrica.
  \item Se repite $B$ veces (t\'ipicamente $B = 1\,000$).
  \item El IC 95\,\% son los percentiles 2,5 y 97,5 de esos $B$ valores.
\end{enumerate}

En el TFM: $B = 1\,000$ bootstraps sobre las predicciones LODO. AUC
pooled = 0,943 con IC 95\,\% [0,925; 0,959]. En TCGA: AUC 0,857 con IC
[0,812; 0,899].

% ========================================================
\section{Aprendizaje autom\'atico: fundamentos}

\subsection{Aprendizaje supervisado}

\begin{itemize}[leftmargin=1.2em]
  \item \textbf{Regresi\'on}: predecir una variable continua (ejemplo:
  supervivencia en meses).
  \item \textbf{Clasificaci\'on}: predecir una etiqueta discreta
  (ejemplo: tumor vs sano, ADC vs SQC).
\end{itemize}

En el TFM: dos problemas de clasificaci\'on binaria.

\subsection{Sobreajuste y sesgo-varianza}

El dilema fundamental del ML:

\begin{itemize}[leftmargin=1.2em]
  \item \textbf{Sesgo alto (underfitting)}: modelo demasiado simple, no
  captura la se\~nal.
  \item \textbf{Varianza alta (overfitting)}: modelo demasiado complejo,
  memoriza el ruido del entrenamiento y falla en datos nuevos.
\end{itemize}

En expresi\'on g\'enica ($p \gg n$, con $p = 20\,000$ genes y $n = 100$
muestras) el riesgo de sobreajuste es extremo. Toda la metodolog\'ia
del TFM (regularizaci\'on L1, validaci\'on LODO externa, consenso
multi-cohorte) est\'a dise\~nada para \emph{controlar la varianza}.

\subsection{Regularizaci\'on}

Modificaci\'on de la funci\'on de coste para penalizar modelos complejos:

\begin{itemize}[leftmargin=1.2em]
  \item \textbf{L2 (Ridge)}: penaliza suma de coeficientes al cuadrado.
  Suaviza. Todos los coeficientes se reducen pero ninguno llega a cero.
  \item \textbf{L1 (LASSO)}: penaliza suma de valores absolutos. Fuerza
  a que muchos coeficientes sean \textbf{exactamente cero}: selecci\'on
  de variables autom\'atica.
  \item \textbf{Elastic Net}: combinaci\'on L1 + L2.
\end{itemize}

\begin{formula}
Funci\'on de coste LASSO para regresi\'on log\'istica:
\[ \mathcal{L} = -\sum_i \log P(y_i|x_i, \beta) + \lambda \sum_j |\beta_j| \]
$\lambda$ (o su inverso $C = 1/\lambda$) es el hiperpar\'ametro que
controla la fuerza de la penalizaci\'on. Se ajusta por validaci\'on
cruzada.
\end{formula}

\subsubsection*{Interpretaci\'on geom\'etrica L1 vs L2}

L1 y L2 son b\'asicamente formas geom\'etricas distintas de la regi\'on
factible:

\begin{itemize}[leftmargin=1.2em]
  \item \textbf{L2 (Ridge)}: regi\'on factible es una \emph{esfera}
  ($\sum \beta_j^2 \leq t$). El \'optimo suele tocarla en un punto donde
  todos los $\beta_j$ son peque\~nos pero no cero.
  \item \textbf{L1 (LASSO)}: regi\'on factible es un \emph{rombo}
  ($\sum |\beta_j| \leq t$). Las esquinas del rombo est\'an sobre los
  ejes coordenados; el \'optimo tiende a tocar en una esquina,
  \textbf{poniendo coeficientes a exactamente cero}.
\end{itemize}

Por eso L1 hace selecci\'on de variables autom\'atica y L2 no.

\subsubsection*{Elecci\'on de $C$: validaci\'on cruzada anidada}

El hiperpar\'ametro $C$ controla el tama\~no del panel:

\begin{itemize}[leftmargin=1.2em]
  \item $C$ peque\~no $\to$ mucha regularizaci\'on $\to$ pocos genes
  con coeficiente no-cero (panel muy peque\~no, quiz\'as sesgado).
  \item $C$ grande $\to$ poca regularizaci\'on $\to$ muchos genes
  seleccionados (panel grande, riesgo de sobreajuste).
\end{itemize}

Se ajusta con \textbf{CV anidada}:

\begin{enumerate}[leftmargin=1.4em]
  \item Bucle externo: LODO (cada cohorte fuera una vez).
  \item Bucle interno: sobre las cohortes de entrenamiento, CV k-fold
  para elegir el mejor $C$.
  \item El $C$ elegido en el bucle interno se aplica al modelo final
  del bucle externo.
\end{enumerate}

La CV anidada evita usar la cohorte de test para elegir el
hiperpar\'ametro (data leakage sutil).

\subsubsection*{Algoritmo de resoluci\'on: coordinate descent}

El t\'ermino $|\beta_j|$ no es diferenciable en cero, por lo que no se
puede resolver con Newton-Raphson est\'andar. La biblioteca
\texttt{scikit-learn} usa \textbf{coordinate descent}: optimiza un
coeficiente por vez manteniendo los dem\'as fijos, aprovechando que
para un solo coeficiente el subproblema tiene soluci\'on cerrada
(operador de \emph{soft-thresholding}).

En el TFM: LASSO L1 sobre regresi\'on log\'istica, con $C$ ajustado por
CV interna. Produce un panel de 20 genes con coeficientes distintos de
cero. Los otros ${\sim}1\,150$ genes candidatos se descartan
autom\'aticamente.

\subsubsection*{Interpretar los coeficientes LASSO}

Los coeficientes $\beta_j$ tras entrenar est\'an en escala
\emph{estandarizada} (todas las variables escaladas a media 0 y
desviaci\'on 1). Un coeficiente $\beta_j = +0{,}45$ significa: por cada
desviaci\'on t\'ipica que aumenta la expresi\'on del gen $j$, el
log-odds de la clase positiva aumenta en 0,45.

\textbf{Cuidado}: en LASSO los coeficientes est\'an sesgados hacia cero
por la penalizaci\'on. Para interpretaci\'on cuantitativa fina se
recomienda \emph{post-selection inference} (reentrenar OLS sobre las
variables seleccionadas) o intervalos de confianza por bootstrap sobre
los coeficientes.

\subsection{Random Forest}

Ensamble de \textbf{\'arboles de decisi\'on} construidos con dos fuentes
de aleatoriedad:

\begin{enumerate}[leftmargin=1.4em]
  \item \textbf{Bagging}: cada \'arbol se entrena sobre una muestra
  bootstrap de los datos (subconjunto con reemplazo).
  \item \textbf{Selecci\'on aleatoria de variables}: en cada divisi\'on
  interna del \'arbol, solo se consideran $\sqrt{p}$ variables aleatorias.
\end{enumerate}

Predicci\'on final: voto mayoritario (clasificaci\'on) o promedio
(regresi\'on). Random Forest ofrece \textbf{importancia de variable}
(cu\'antas veces se selecciona cada gen en las divisiones) como
complemento a la interpretaci\'on.

\subsubsection*{Hiperpar\'ametros clave}

\begin{itemize}[leftmargin=1.2em]
  \item \textbf{\texttt{n\_estimators}}: n\'umero de \'arboles.
  M\'as = mejor pero con retorno decreciente. 100-500 t\'ipico.
  \item \textbf{\texttt{max\_depth}}: profundidad m\'axima de cada
  \'arbol. Si \texttt{None}, crecen hasta hojas puras (overfitting en
  algunos casos).
  \item \textbf{\texttt{min\_samples\_leaf}}: m\'inimo de muestras por
  hoja. Regularizaci\'on efectiva.
  \item \textbf{\texttt{max\_features}}: variables consideradas en cada
  divisi\'on. Por defecto $\sqrt{p}$ en clasificaci\'on.
\end{itemize}

\subsubsection*{Error OOB (out-of-bag)}

Cada \'arbol se entrena con una muestra bootstrap ($\sim$63\,\% de las
muestras originales). El $\sim$37\,\% restante son \emph{out-of-bag} para
ese \'arbol. Se puede evaluar el error del bosque promediando
predicciones OOB. \textbf{Es un estimador insesgado del error de
generalizaci\'on sin necesidad de CV}.

\subsubsection*{Importancia de variable}

Dos m\'etodos principales:
\begin{itemize}[leftmargin=1.2em]
  \item \textbf{Mean Decrease Impurity (MDI, Gini importance)}:
  cu\'anto reduce la impureza cada variable al usarse en divisiones. R\'apida
  pero \emph{sesgada hacia variables continuas o con muchas categor\'ias}.
  \item \textbf{Permutation importance}: permutar valores de una variable
  y medir la ca\'ida de rendimiento. M\'as costosa pero no sesgada.
\end{itemize}

En el TFM: Random Forest se usa como \emph{sanity check} vs LASSO,
confirmando que los genes seleccionados por LASSO tambi\'en aparecen
altos en el ranking RF.

\subsection{SVM lineal}

Support Vector Machine busca el \textbf{hiperplano \'optimo} que separa
dos clases maximizando el margen (distancia a los puntos m\'as cercanos
de cada clase). En su versi\'on lineal es esencialmente una regresi\'on
log\'istica con otra funci\'on de p\'erdida (\emph{hinge loss}) y otra
regularizaci\'on. En dimensiones altas rinde bien, pero los pesos no
son directamente interpretables como ``importancia de gen''.

\begin{formula}
\textbf{Formulaci\'on primal SVM lineal (soft-margin)}:
\[ \min_{w, b, \xi} \quad \frac{1}{2}\|w\|^2 + C\sum_i \xi_i \]
sujeto a $y_i(w \cdot x_i + b) \geq 1 - \xi_i$, $\xi_i \geq 0$.

$C$ controla la tolerancia a errores: $C$ grande $\to$ margen estrecho,
poca tolerancia; $C$ peque\~no $\to$ margen amplio, m\'as tolerancia.
Los puntos con $\xi_i > 0$ son \emph{vectores soporte}: definen el
hiperplano.
\end{formula}

\subsubsection*{Kernel trick}

Para relaciones no lineales, SVM proyecta impl\'icitamente los datos a
un espacio de mayor dimensi\'on donde s\'i son linealmente separables,
usando funciones kernel $K(x_i, x_j)$:

\begin{itemize}[leftmargin=1.2em]
  \item \textbf{Lineal}: $K(x,x') = x \cdot x'$. Es SVM lineal.
  \item \textbf{Polinomial}: $K(x,x') = (\gamma \, x \cdot x' + r)^d$.
  \item \textbf{RBF (Gaussiano)}: $K(x,x') = \exp(-\gamma \|x-x'\|^2)$.
  El m\'as popular.
\end{itemize}

En el TFM se usa \emph{s\'olo} SVM lineal por interpretabilidad
comparativa (kernel no lineal complica la extracci\'on de importancia
de gen).

\begin{concepto}{Por qu\'e LASSO y no SVM en el TFM}
Comparativa LODO: LASSO 0,925, RF 0,933, SVM 0,957 (AUC media). SVM
gana en rendimiento por 3 puntos, pero: (i) los coeficientes SVM no son
directamente interpretables cl\'inicamente, (ii) LASSO ya produce un
panel discreto de 20 genes traducible a NanoString / RT-qPCR, (iii) la
diferencia de rendimiento es menor que la variabilidad entre cohortes.
LASSO se elige por su balance \textbf{rendimiento-interpretabilidad}.
\end{concepto}

% ========================================================
\section{Validaci\'on cruzada, LODO y por qu\'e importa}

\subsection{Validaci\'on cruzada est\'andar (k-fold)}

Se divide el conjunto de datos en $k$ partes iguales (folds). Se entrena
con $k-1$ y se eval\'ua con la restante. Se rota $k$ veces. Al final se
promedia el rendimiento. $k = 5$ o $k = 10$ son comunes.

Problema en gen\'omica multi-cohorte: si mezclas muestras de todas las
cohortes en los folds, dos muestras del mismo estudio pueden acabar una
en train y otra en test. Como comparten sesgos t\'ecnicos (chip, lote,
laboratorio), el modelo \emph{aprende} el sesgo t\'ecnico y da AUCs
inflados.

\subsection{Leave-One-Dataset-Out (LODO)}

Variante estricta: cada \emph{cohorte completa} se retira una vez del
entrenamiento y se usa como test. Si hay 8 cohortes:

\begin{itemize}[leftmargin=1.2em]
  \item Iteraci\'on 1: entrenar con 7, testear con cohorte 1.
  \item Iteraci\'on 2: entrenar con 7 (otra combinaci\'on), testear con 2.
  \item \dots\ 8 iteraciones en total.
\end{itemize}

Ninguna muestra de la cohorte de test aparece nunca en el entrenamiento.
Es el est\'andar oro de validaci\'on externa en meta-an\'alisis
transcript\'omicos. El AUC bajo LODO es t\'ipicamente 5-15 puntos m\'as
bajo que bajo CV est\'andar --- porque LODO \emph{no infla} por sesgo
compartido de cohorte.

\subsubsection*{Otras estrategias de partici\'on}

\begin{itemize}[leftmargin=1.2em]
  \item \textbf{Hold-out}: dividir 80/20. R\'apido pero de alta varianza.
  \item \textbf{Stratified k-fold}: mantiene la proporci\'on de clases
  en cada fold.
  \item \textbf{Group k-fold}: mantiene grupos (paciente, familia) sin
  cruzar entre folds.
  \item \textbf{Leave-One-Out (LOO)}: caso extremo con $k = n$. Muy caro
  y de alta varianza en gen\'omica.
  \item \textbf{Nested CV}: CV externa + CV interna para elegir
  hiperpar\'ametros. Est\'andar cuando hay hiperpar\'ametros a ajustar.
\end{itemize}

En el TFM se combina \textbf{LODO externa} (cohorte fuera) con
\textbf{k-fold interna} sobre las cohortes de train para elegir el
$C$ del LASSO. Es CV anidada donde la partici\'on externa respeta grupos
de cohorte.

\begin{concepto}{Frase para la defensa}
\emph{``LODO no corrige el efecto lote entre cohortes: lo mide. Si la
firma resulta robusta bajo LODO, es porque los genes seleccionados
mantienen su patr\'on de cambio a pesar del efecto lote, lo que es una
propiedad deseable.''}
\end{concepto}

\subsection{Batch effect (efecto lote)}

Diferencias sistem\'aticas entre muestras atribuibles al lote t\'ecnico
(chip usado, reactivos, operador, fecha, laboratorio). En cohortes GEO
combinadas, gran parte de la varianza total viene del efecto lote y no
de la biolog\'ia.

\subsubsection*{ComBat y por qu\'e NO lo usamos}
ComBat (Johnson et al., 2007) es un m\'etodo bayesiano emp\'irico que
``corrige'' el efecto lote normalizando las cohortes conjuntamente.
Problema: si el efecto lote est\'a \emph{confundido} con la variable de
inter\'es (ejemplo: un centro solo aporta tumores, otro solo controles),
ComBat borra tambi\'en la se\~nal biol\'ogica.

En el TFM: no aplicamos ComBat. En su lugar (i) normalizamos
\emph{dentro} de cada cohorte, (ii) exigimos consenso multi-cohorte
para seleccionar genes, (iii) usamos LODO como criterio de rendimiento.
Estrategia conservadora pero honesta.

% ========================================================
\section{M\'etricas y calibraci\'on}

\subsection{Matriz de confusi\'on y m\'etricas base}

Para clasificaci\'on binaria (positivo/negativo):

\begin{center}
\small
\begin{tabular}{lcc}
\toprule
 & \textbf{Predicho +} & \textbf{Predicho $-$} \\
\midrule
\textbf{Real +} & TP (verdadero positivo) & FN (falso negativo) \\
\textbf{Real $-$} & FP (falso positivo) & TN (verdadero negativo) \\
\bottomrule
\end{tabular}
\end{center}

M\'etricas:

\begin{itemize}[leftmargin=1.2em]
  \item \textbf{Sensibilidad (recall, TPR)}: $\text{TP} / (\text{TP} + \text{FN})$.
  De todos los positivos reales, cu\'antos detecto.
  \item \textbf{Especificidad (TNR)}: $\text{TN} / (\text{TN} + \text{FP})$.
  De todos los negativos reales, cu\'antos descarto correctamente.
  \item \textbf{Precisi\'on (PPV)}: $\text{TP} / (\text{TP} + \text{FP})$.
  De los que predije positivos, cu\'antos lo son. \textbf{Depende de la
  prevalencia}.
  \item \textbf{Balanced Accuracy}: media de sensibilidad y especificidad.
  \'Util con clases desbalanceadas.
  \item \textbf{Accuracy ingenua}: $(\text{TP}+\text{TN})/N$. En clases
  desbalanceadas puede ser enga\~nosa (predecir siempre la clase
  mayoritaria puede dar 90\,\%).
\end{itemize}

\subsection{Curva ROC y AUC}

\textbf{ROC} (Receiver Operating Characteristic): representa sensibilidad
frente a $1-$especificidad para todos los umbrales de decisi\'on. El
\'area bajo esta curva (\textbf{AUC}) es una m\'etrica \emph{umbral-independiente}:

\begin{itemize}[leftmargin=1.2em]
  \item AUC = 0,5: modelo al azar.
  \item AUC = 1,0: perfecto.
  \item AUC $\geq 0{,}9$: excelente discriminaci\'on.
  \item AUC $\geq 0{,}8$: buena discriminaci\'on.
\end{itemize}

Interpretaci\'on probabil\'istica: AUC es la probabilidad de que el
modelo ordene una muestra positiva por encima de una negativa. \'Util
para comparar modelos independientemente del umbral, pero \emph{no dice}
si las probabilidades est\'an calibradas.

\subsubsection*{Curva Precision-Recall (PR)}

Alternativa a ROC especialmente \'util en \textbf{clases muy
desbalanceadas}. Representa precisi\'on vs recall.

\begin{itemize}[leftmargin=1.2em]
  \item ROC puede parecer ``buena'' incluso con clasificador in\'util si
  hay mucho desbalance (TNR alto por defecto).
  \item PR es m\'as informativa cuando la clase positiva es rara.
\end{itemize}

En el TFM: las clases est\'an moderadamente balanceadas en las cohortes
evaluables (ratio tumor:sano suele estar entre 1:1 y 3:1), por lo que
ROC es adecuada. Cuando la desviaci\'on es extrema (GSE31210 tiene
226 tumores vs 20 sanos), reportamos tambi\'en balanced accuracy y
sensibilidad/especificidad por separado.

\subsubsection*{DeLong test para comparar AUCs}

Para comparar dos AUCs sobre el mismo conjunto de test (por ejemplo LASSO
vs SVM), no se puede simplemente comparar los valores puntuales. El test
de DeLong (1988) da un $p$-valor para $H_0: AUC_1 = AUC_2$ teniendo en
cuenta la correlaci\'on entre los rankings.

En el TFM no aplicamos DeLong formalmente porque nos apoyamos en los
IC bootstrap para comparar modelos.

\subsection{Calibraci\'on de probabilidades}

Un modelo con AUC alto puede tener probabilidades \textbf{descalibradas}:
predice 0,9 y en realidad la clase real ocurre el 60\,\% de las veces.
La calibraci\'on es distinta de la discriminaci\'on:

\begin{itemize}[leftmargin=1.2em]
  \item \textbf{Platt scaling}: ajusta una regresi\'on log\'istica sobre
  la puntuaci\'on cruda del modelo. Asume forma sigmoidea.
  \item \textbf{Regresi\'on isot\'onica}: aprende una funci\'on
  \emph{mon\'otona no param\'etrica}. M\'as flexible pero requiere m\'as
  datos.
\end{itemize}

En el TFM: se aplican ambas sobre las predicciones LODO. Isot\'onica
mejora BalAcc de 0,772 $\to$ 0,810 y AUC de 0,925 $\to$ 0,953. Platt
degrada el modelo (BalAcc 0,772 $\to$ 0,639). Por eso se selecciona
isot\'onica.

\begin{formula}
\textbf{Brier score}: m\'etrica de calibraci\'on. Media del cuadrado
del error entre probabilidad predicha $p_i$ y etiqueta real $y_i \in
\{0,1\}$:
\[ BS = \frac{1}{n}\sum_i (p_i - y_i)^2 \]
Menor es mejor. $BS = 0$: perfecta calibraci\'on y discriminaci\'on.
\end{formula}

\subsubsection*{Diagrama de fiabilidad (reliability plot)}

Herramienta visual de calibraci\'on:

\begin{enumerate}[leftmargin=1.4em]
  \item Se dividen las muestras en bins seg\'un su probabilidad
  predicha (ejemplo: 10 bins de ancho 0,1).
  \item Para cada bin se calcula la \emph{frecuencia observada} de la
  clase positiva.
  \item Se representa probabilidad predicha (x) vs frecuencia observada
  (y).
  \item Un modelo perfectamente calibrado cae sobre la diagonal $y = x$.
\end{enumerate}

Otras m\'etricas: ECE (Expected Calibration Error), MCE (Maximum
Calibration Error).

\subsubsection*{Por qu\'e isot\'onica gan\'o a Platt en el TFM}

\begin{itemize}[leftmargin=1.2em]
  \item Platt scaling asume que la relaci\'on score-probabilidad es
  sigmoidea. Si no lo es, no puede corregirla.
  \item Isot\'onica es no param\'etrica: solo asume monotonicidad. En
  cohortes heterog\'eneas, la relaci\'on real puede no ser sigmoidea y
  isot\'onica la captura mejor.
  \item Riesgo de isot\'onica: puede sobreajustar con pocas muestras.
  Con $n < 100$ Platt suele ser preferible.
\end{itemize}

En el TFM los conjuntos LODO por cohorte son suficientes ($n \geq 100$
en la mayor\'ia) para que isot\'onica sea robusta.

\subsubsection*{Manejo del desbalance de clases}

Cuando las clases est\'an muy desbalanceadas ($>1:5$), el modelo tiende
a predecir la clase mayoritaria. Estrategias:

\begin{itemize}[leftmargin=1.2em]
  \item \textbf{Pesos de clase}: multiplicar la p\'erdida por
  $w_c = n / (K \cdot n_c)$. \textbf{Estrategia por defecto en el TFM}
  (\texttt{class\_weight="balanced"} en scikit-learn).
  \item \textbf{Oversampling}: duplicar la clase minoritaria (ejemplo:
  SMOTE, que genera muestras sint\'eticas por interpolaci\'on).
  \item \textbf{Undersampling}: submuestrear la clase mayoritaria. Pierde
  informaci\'on.
  \item \textbf{Ajustar el umbral de decisi\'on}: si el AUC es bueno pero
  el modelo predice todo como mayoritario, bajar el umbral.
\end{itemize}

% ========================================================
\section{Decisiones de dise\~no del pipeline}

Anticipaci\'on a preguntas del tribunal ``por qu\'e X y no Y''.

\subsection{Por qu\'e microarrays y no s\'olo RNA-Seq}

\begin{itemize}[leftmargin=1.2em]
  \item Hay muchas m\'as cohortes de microarray en GEO con metadatos
  ricos para ADC vs SQC.
  \item RNA-Seq es plataforma m\'as reciente; combinar decenas de
  cohortes RNA-Seq con metadatos consistentes es dif\'icil.
  \item La firma se valida despu\'es en RNA-Seq (TCGA), demostrando
  transferibilidad.
\end{itemize}

\subsection{Por qu\'e consenso multi-cohorte y no ComBat}

Cubierto arriba: ComBat puede borrar se\~nal biol\'ogica cuando el lote
confunde con la variable de inter\'es. Consenso es m\'as conservador.

\subsection{Por qu\'e $|d| > 0{,}5$ y no otro umbral}

Umbral cl\'asico de Cohen para efecto ``moderado''. Es un compromiso:
\begin{itemize}[leftmargin=1.2em]
  \item $|d| > 0{,}2$: demasiado permisivo, entran genes con efecto
  peque\~no que pueden ser artefacto.
  \item $|d| > 0{,}8$: demasiado exigente, se pierden genes con efecto
  moderado pero muy reproducible.
  \item $|d| > 0{,}5$: convenci\'on est\'andar, filtra ruido sin
  descartar biolog\'ia leg\'itima.
\end{itemize}

\subsection{Por qu\'e 3 cohortes de consenso y no 5 o 8}

Trade-off entre reproducibilidad y disponibilidad:

\begin{itemize}[leftmargin=1.2em]
  \item Exigir 8 cohortes reduce dram\'aticamente el n\'umero de genes
  que pasan y perder\'ia se\~nal cl\'inicamente relevante.
  \item Exigir 3 cohortes independientes ya es m\'as estricto que la
  mayor\'ia de firmas publicadas (que validan en 0 o 1 cohorte
  independiente).
  \item El resto de cohortes (5) se usan como test LODO, no como
  entrenamiento, para validaci\'on externa adicional.
\end{itemize}

\subsection{Por qu\'e panel m\'inimo de 20 genes y no 5 o 50}

Curva panel vs AUC (\texttt{PANEL\_MINIMO\_CURVA.csv}):

\begin{itemize}[leftmargin=1.2em]
  \item Con 5 genes: AUC $\approx 0{,}91$. Muy compacto pero pierde
  redundancia biol\'ogica.
  \item Con 20 genes: AUC $\approx 0{,}966$. Plateau: a\~nadir m\'as no
  mejora significativamente.
  \item Con 50 genes: AUC $\approx 0{,}97$. Rendimiento pr\'acticamente
  id\'entico al panel completo.
\end{itemize}

20 es el ``codo'' de la curva y coincide con el tama\~no del panel OMS
IHC, lo que facilita la interpretaci\'on comparativa.

\subsection{Por qu\'e recalibraci\'on isot\'onica y no Platt}

Cubierto arriba: isot\'onica gan\'o experimentalmente en el TFM
(BalAcc 0,772 $\to$ 0,810 vs Platt 0,772 $\to$ 0,639).

\subsection{Por qu\'e bootstrap y no otro m\'etodo de IC}

\begin{itemize}[leftmargin=1.2em]
  \item \textbf{IC anal\'iticos} (asumen normalidad): AUC no tiene una
  distribuci\'on cerrada simple.
  \item \textbf{Jackknife}: variante de bootstrap con leave-one-out.
  M\'as barato pero de mayor sesgo.
  \item \textbf{Bootstrap percentil}: est\'andar oro para m\'etricas
  no lineales como AUC. \textbf{Elegido en el TFM}.
\end{itemize}

% ========================================================
\section{Modelos de lenguaje aplicados a bioinform\'atica}

\subsection{Qu\'e es un LLM}

Un modelo de lenguaje grande (LLM) es una red neuronal profunda
entrenada con enormes cantidades de texto para predecir la siguiente
palabra dado un contexto. Arquitectura dominante: \textbf{transformer}
(Vaswani et al., 2017), basado en el mecanismo de \emph{atenci\'on}.

Escalas modernas: 1--700 mil millones de par\'ametros. Modelos
representativos: GPT-4 (OpenAI, cerrado), Claude (Anthropic, cerrado),
Gemini (Google), Llama (Meta, abierto), Mistral (europeo, abierto).

\subsection{Llama 3.3-70b}

Publicado por Meta AI a finales de 2024. 70 mil millones de par\'ametros.
Rendimiento comparable a GPT-4o en benchmarks acad\'emicos. Ventaja
clave: es \textbf{open weights}, se puede desplegar donde uno quiera,
sin depender de una API cerrada ni exponer datos sensibles.

\subsection{Groq}

Empresa que desarrolla hardware espec\'ifico (\textbf{LPU},
\emph{Language Processing Unit}) optimizado para inferencia de LLMs.
Consigue latencias 5--10$\times$ m\'as r\'apidas que GPU tradicionales
para el mismo modelo. Ofrece API p\'ublica que sirve Llama con costes
muy bajos y velocidad de decenas de tokens por segundo.

\subsection{Rol del LLM en el TFM}

Tarea \'unica y bien delimitada: leer una descripci\'on de muestra GEO
(texto libre) y devolver una etiqueta can\'onica del grupo experimental.

\textbf{Prompt estructurado}:
\begin{itemize}[leftmargin=1.2em]
  \item Rol: ``eres un experto anotador de datos biom\'edicos''.
  \item Contexto: descripci\'on de la muestra + metadatos del estudio.
  \item Instrucci\'on: elige entre lista cerrada de etiquetas.
  \item Regla de escape: si no encaja ninguna, devuelve \texttt{unknown}.
\end{itemize}

Tasa de \'exito global: 82,8\,\% (curados sobre descargados). Las
muestras que el LLM no puede clasificar se descartan, no se fuerzan.

\subsection{Alucinaciones y c\'omo se controlan}

\emph{Alucinaci\'on}: el LLM inventa informaci\'on que no est\'a en el
prompt. Es un modo de fallo conocido.

Mitigaciones en el TFM:
\begin{itemize}[leftmargin=1.2em]
  \item Tarea de clasificaci\'on cerrada, no generaci\'on abierta.
  \item Lista de etiquetas explicitada; ``unknown'' permitido.
  \item El LLM no participa en el an\'alisis estad\'istico ni en la
  selecci\'on de genes.
  \item Auditor\'ia manual sobre un subconjunto para verificar la
  precisi\'on real.
\end{itemize}

\begin{peligro}
El uso de LLMs en bioinform\'atica todav\'ia levanta escepticismo. Si te
preguntan ``no le est\'as dando poder al modelo?'' responde con claridad:
el LLM \textbf{s\'olo cura metadatos}. Todo el an\'alisis estad\'istico,
la selecci\'on de genes, la clasificaci\'on y la validaci\'on son
determin\'isticos, reproducibles con \texttt{random\_state=42} y no
dependen del LLM.
\end{peligro}

% ========================================================
\section{Consenso multi-cohorte y meta-an\'alisis}

\subsection{Por qu\'e integrar cohortes}

Un solo estudio tiene poder estad\'istico limitado, sesgos particulares y
puede detectar se\~nales que no replican. Integrar varios estudios
aumenta el tama\~no muestral efectivo y permite distinguir se\~nales
reproducibles de artefactos.

Dos estrategias generales:
\begin{itemize}[leftmargin=1.2em]
  \item \textbf{Combinaci\'on con correcci\'on de lote} (ejemplo: ComBat).
  Simple pero arriesgado si el lote confunde con la variable de inter\'es.
  \item \textbf{Meta-an\'alisis por consenso}. Cada cohorte se analiza
  por separado; se combinan los resultados con criterios estad\'isticos.
\end{itemize}

El TFM adopta la segunda.

\subsection{Consenso en el TFM}

Para cada gen $g$ y cada cohorte $c$:
\begin{itemize}[leftmargin=1.2em]
  \item Se calcula el $d$ de Cohen entre grupos (ejemplo: ADC vs SQC).
  \item Se registra el signo y la magnitud.
\end{itemize}

Un gen \emph{sobrevive} al consenso si en \textbf{las 3 cohortes
independientes}:
\begin{itemize}[leftmargin=1.2em]
  \item Mantiene el \emph{mismo signo} (misma direcci\'on de cambio).
  \item Alcanza $|d| > 0{,}5$ (efecto moderado o mayor).
\end{itemize}

Resultado: de $\sim$22\,880 genes evaluados, sobreviven \textbf{1\,174}.
Es un filtro exigente que garantiza reproducibilidad biol\'ogica.

% ========================================================
\section{Interpretabilidad: ejes biol\'ogicos}

Una firma de 1\,174 genes puede parecer una caja negra. En el TFM se
analiza la estructura interna con dos preguntas:

\begin{enumerate}[leftmargin=1.4em]
  \item \textbf{Los genes seleccionados coinciden con marcadores
  cl\'inicos ya conocidos?} S\'i: 18 de 20 marcadores IHC del panel OMS
  aparecen en la firma sin haberse declarado.
  \item \textbf{Existen ejes biol\'ogicos coherentes en la firma?} S\'i,
  dos:
\end{enumerate}

\subsubsection*{Eje 1: p\'erdida alv\'eolo-capilar}
Genes como \gene{SFTPC}, \gene{SFTPB}, \gene{AGER}, \gene{CLDN18},
\gene{FABP4}, marcadores canonicos de neumocitos y capilares alveolares.
En muestras tumorales la expresi\'on de estos marcadores cae. La firma
correlaciona negativamente con estos marcadores ($\rho = -0{,}626$ media
en 4 cohortes con muestra suficiente). Biol\'ogicamente coherente: el
tumor sustituye al par\'enquima alveolar sano.

\subsubsection*{Eje 2: proliferaci\'on tumoral}
Genes como \gene{MKI67}, \gene{TOP2A}, \gene{CCNB1}, \gene{BIRC5},
\gene{AURKA}. Marcadores cl\'asicos de proliferaci\'on que el pat\'ologo
usa en cl\'inica (Ki-67 en biopsias). Un score alto en este eje
correlaciona con tumores m\'as agresivos.

\begin{concepto}{Frase para el tribunal}
\emph{``La firma no es una caja negra. Captura la transici\'on biol\'ogica
del pulm\'on sano al tejido tumoral en dos ejes que un pat\'ologo
reconoce: p\'erdida de arquitectura alveolar y ganancia de actividad
proliferativa.''}
\end{concepto}

% ========================================================
\section{Transferibilidad cl\'inica}

\subsection{NanoString nCounter}

Tecnolog\'ia que mide expresi\'on de un panel peque\~no de genes
(decenas a cientos) directamente sobre tejido fijado en formalina
(\textbf{FFPE}). El tejido FFPE es el est\'andar cl\'inico en anatom\'ia
patol\'ogica (se conserva a\~nos, es barato). NanoString no requiere
RT-PCR y da lecturas cuantitativas robustas. Ideal para paneles
diagn\'osticos multiplexados.

\subsection{RT-qPCR}

PCR cuantitativa a tiempo real. Est\'andar cl\'inico para paneles muy
peque\~nos (1-20 genes). Alta sensibilidad, coste bajo, resultado en
horas. Requiere ARN de buena calidad (m\'as exigente que NanoString).

\subsection{Por qu\'e importa al TFM}

Los 20 genes del panel m\'inimo son \emph{traducibles} a estas
plataformas. Podr\'ian construirse ensayos NanoString / qPCR para uso
diagn\'ostico. Diferencia clave respecto a firmas basadas en RNA-Seq
completo, que no son viables en flujo cl\'inico rutinario.

% ========================================================
\section{Reproducibilidad}

\subsection{Principios}

Una investigaci\'on reproducible cumple: cualquier persona con acceso
al c\'odigo y a los datos originales puede regenerar todas las cifras y
figuras publicadas.

Elementos del TFM que garantizan reproducibilidad:

\begin{itemize}[leftmargin=1.2em]
  \item \textbf{Semilla aleatoria fija}: \texttt{random\_state=42} en
  todos los pasos aleatorizados.
  \item \textbf{C\'odigo p\'ublico}: repositorio Git p\'ublico con
  historial completo.
  \item \textbf{Datos p\'ublicos}: cohortes GEO y TCGA descargables por
  cualquiera.
  \item \textbf{Artefactos intermedios versionados}: cada paso escribe
  CSV/JSON que se pueden auditar.
  \item \textbf{Especificaci\'on de entorno}: \texttt{requirements.txt}
  con versiones exactas.
\end{itemize}

\subsection{Data leakage}

\emph{Data leakage} (fuga de datos): informaci\'on del conjunto de test
se filtra al entrenamiento. Es el error m\'as com\'un en ML aplicado a
gen\'omica.

Fuentes t\'ipicas:
\begin{itemize}[leftmargin=1.2em]
  \item Selecci\'on de features en todo el conjunto antes de dividir.
  \item Normalizaci\'on aplicada al conjunto completo.
  \item Uso de la misma cohorte en train y test (por descuido).
\end{itemize}

En el TFM: la selecci\'on de genes y la calibraci\'on se hacen
\emph{dentro} de cada iteraci\'on LODO, usando s\'olo las cohortes de
entrenamiento. Ninguna muestra de la cohorte test se ve nunca durante
el ajuste.

En TCGA: las 408 muestras TCGA se aplican \emph{despu\'es} de entrenar
todo el pipeline en microarray. Ninguna muestra TCGA particip\'o en la
selecci\'on de features ni en la calibraci\'on.

% ========================================================
\newpage
\section{Glosario r\'apido}

\begin{description}[leftmargin=0em, style=nextline]
  \item[ADC] Adenocarcinoma. Subtipo de NSCLC de origen glandular.
  \item[ARN mensajero (mRNA)] Copia funcional de un gen, sustrato de la traducci\'on.
  \item[AUC] Area Under the ROC Curve. Discriminaci\'on de un clasificador.
  \item[Balanced accuracy] Media de sensibilidad y especificidad.
  \item[Batch effect] Diferencias t\'ecnicas sistem\'aticas entre lotes.
  \item[Bootstrap] Remuestreo con reemplazo para estimar incertidumbre.
  \item[cBioPortal] Portal p\'ublico que sirve datos TCGA v\'ia API.
  \item[Cohen ($d$)] Tama\~no de efecto estandarizado.
  \item[ComBat] Algoritmo bayesiano de correcci\'on de efecto lote.
  \item[Data leakage] Contaminaci\'on inadvertida train-test.
  \item[$d$ de Cohen] Ver ``Cohen''.
  \item[FDR] False Discovery Rate. Correcci\'on por m\'ultiples tests (BH).
  \item[FFPE] Tejido fijado en formalina y embebido en parafina. Est\'andar cl\'inico.
  \item[Firma g\'enica] Conjunto de genes cuya expresi\'on caracteriza un fenotipo.
  \item[GEO] Gene Expression Omnibus. Repositorio p\'ublico del NCBI.
  \item[GPL] Identificador de plataforma tecnol\'ogica en GEO.
  \item[Groq] Empresa de inferencia LLM sobre hardware LPU.
  \item[GSE] Identificador de estudio en GEO.
  \item[GSM] Identificador de muestra individual en GEO.
  \item[HGNC] Human Genome Organisation. Nomenclatura g\'enica oficial.
  \item[IHC] Inmunohistoqu\'imica. T\'ecnica anatomopatol\'ogica.
  \item[Isot\'onica (regresi\'on)] M\'etodo de calibraci\'on no param\'etrico.
  \item[Ki-67 (MKI67)] Marcador cl\'asico de proliferaci\'on.
  \item[k-fold CV] Validaci\'on cruzada en $k$ particiones.
  \item[LASSO] Regresi\'on con penalizaci\'on L1. Selecciona variables.
  \item[Llama] Familia de LLMs abiertos de Meta.
  \item[LODO] Leave-One-Dataset-Out. Validaci\'on externa por cohorte.
  \item[LPU] Hardware espec\'ifico de Groq para inferencia LLM.
  \item[LUAD / LUSC] Nombres TCGA para adenocarcinoma / escamoso de pulm\'on.
  \item[Microarray] Chip que mide expresi\'on por hibridaci\'on.
  \item[NanoString] Plataforma comercial para paneles de genes en FFPE.
  \item[NCBI] National Center for Biotechnology Information (EE.UU.).
  \item[NSCLC] Non-Small Cell Lung Cancer.
  \item[OMS] Organizaci\'on Mundial de la Salud. Publica clasificaci\'on de tumores.
  \item[Panel OMS 2015] 20 marcadores IHC est\'andar para NSCLC (Travis et al., 2015).
  \item[Pipeline] Cadena de pasos automatizados que produce resultados.
  \item[Platt scaling] Calibraci\'on sigmoidea con regresi\'on log\'istica.
  \item[Random Forest] Ensamble de \'arboles de decisi\'on con bagging.
  \item[Reproducibilidad] Capacidad de regenerar todas las cifras publicadas.
  \item[RNA-Seq] Secuenciaci\'on del transcriptoma.
  \item[ROC] Curva sensibilidad vs 1-especificidad.
  \item[RSEM] Estimador estad\'istico para cuantificar RNA-Seq. Usado por TCGA.
  \item[SCLC] Small Cell Lung Cancer. Subtipo minoritario y agresivo.
  \item[Sensibilidad (recall, TPR)] TP / (TP + FN).
  \item[Spearman ($\rho$)] Correlaci\'on por rangos, mide monoton\'ia.
  \item[SQC / LUSC] Squamous Cell Carcinoma (escamoso).
  \item[SVM] Support Vector Machine. Clasificador de m\'aximo margen.
  \item[TCGA] The Cancer Genome Atlas. Consorcio NIH de caracterizaci\'on tumoral.
  \item[TNM] Sistema de estadificaci\'on tumor-ganglios-met\'astasis.
  \item[TPM] Transcripts Per Million. Unidad de normalizaci\'on RNA-Seq.
  \item[Transcriptoma] Conjunto de mRNAs presentes en una muestra.
  \item[Transformer] Arquitectura de red neuronal base de los LLMs.
  \item[TTF-1 (NKX2-1)] Marcador IHC de adenocarcinoma.
  \item[Welch (t de)] Test t para varianzas desiguales.
\end{description}

\newpage

\begin{tcolorbox}[colback=plano, colframe=verdeTerm, arc=3pt, boxrule=1pt]
\centering\itshape
Este manual cubre los fundamentos que sostienen el TFM. Si te preguntan
algo que no est\'a aqu\'i, la respuesta honesta ``no lo s\'e con
precisi\'on, pero mi trabajo se apoya en \dots'' es preferible a
inventar. La memoria est\'a bien construida y los resultados son
reproducibles: no necesitas defender m\'as de lo que ya has hecho.
\end{tcolorbox}

\end{document}
